\documentclass[11pt]{article}
\usepackage[margin=1in]{geometry}
\usepackage{amsmath,amssymb,amsthm}
\usepackage{graphicx}
\usepackage{hyperref}

\newtheorem{theorem}{Theorem}
\newtheorem{remark}{Remark}

\newcommand{\bpi}{\bar\pi}
\newcommand{\dd}{\mathrm{d}}
\newcommand{\im}{\operatorname{im}}
\newcommand{\Isom}{\operatorname{Isom}}
\newcommand{\C}{\mathbb C}
\newcommand{\R}{\mathbb R}

\title{A finite-dimensional case of Nowak's kernel-pairing question}
\author{FA/Banach automatic discovery run}
\date{2026-07-03}

\begin{document}
\maketitle

\noindent\textbf{Status.} Candidate partial result, likely valid. This answers the
finite-dimensional case of Remark 4.3 in Nowak's paper \cite{Nowak2011}; it
does not address infinite-dimensional isometric representations.

\section*{Source question}

The source is Piotr W. Nowak, \emph{Poincare inequalities and rigidity for
actions on Banach spaces}, arXiv:1107.1896 \cite{Nowak2011}. Remark 4.3 asks
for which isometric representations the natural map
\[
   i^*\circ \bar i:\ker d_{\bar\pi}\longrightarrow (\ker d_\pi)^*
\]
is automatically an isomorphism, or at least surjective. The paper notes that
this would eliminate, for such representations, the need for one of the two
Poincare-constant inequalities in the proof of the main theorem.

\begin{center}
\includegraphics[width=\textwidth]{figures/open_problem_crop_page15.png}

\includegraphics[width=\textwidth]{figures/open_problem_crop_page16.png}
\end{center}

\section*{Idea of the proof}

Finite-dimensional isometric representations have compact image. Averaging an
auxiliary Euclidean or Hermitian inner product over the compact closure of the
image makes the representation unitary for that auxiliary inner product. The
Riesz map associated with this invariant inner product intertwines the original
representation with the contragredient representation. Applied coordinatewise
to Nowak's cochain spaces, it carries \(\ker d_\pi\) onto \(\ker d_{\bar\pi}\).
The source pairing then becomes a positive definite Hilbert sum on the graph of
this identification, so no nonzero kernel vector can be annihilated by all dual
kernel vectors. Finite dimensionality converts this separation into
surjectivity, indeed into an isomorphism.

\section*{Definitions}

We use the notation of \cite{Nowak2011}. Let \(G\) be generated by a finite
symmetric set \(S\) equipped with an admissible weight \(\omega\) on the link
graph. Let \(X\) be a finite-dimensional normed space and let
\(\pi:G\to \Isom(X)\) be an isometric representation. The contragredient
representation on \(X^*\) is
\[
   \bar\pi_g=\pi_{g^{-1}}^* .
\]
The cochain space \(C_-^{(1,p)}(G,\pi)\) consists of functions \(f:S\to X\)
with \(f=-Q_\pi f\), and \(d_\pi:C_-^{(1,p)}(G,\pi)\to C^{(2,p)}(G,\pi)\) is
Nowak's operator
\[
   d_\pi f(s,t)=\pi_s f(s^{-1}t)-f(t)+f(s).
\]
The map in Remark 4.3 is the restriction-pairing map
\[
 R_\pi:\ker d_{\bar\pi}\longrightarrow (\ker d_\pi)^*,\qquad
 R_\pi(\varphi)(f)=\sum_{s\in S}\varphi(s)(f(s))\,\deg_\omega(s).
\]

\begin{theorem}
Let \(X\) be a finite-dimensional real or complex normed space, and let
\(\pi:G\to\Isom(X)\) be an isometric representation. Then
\[
   R_\pi:\ker d_{\bar\pi}\longrightarrow(\ker d_\pi)^*
\]
is an isomorphism. In particular, the map \(i^*\circ\bar i\) in Nowak's
Remark 4.3 is automatically surjective for every finite-dimensional isometric
representation.
\end{theorem}

\begin{proof}
The group \(\Isom(X)\) is compact in finite dimension: it is a closed and
bounded subgroup of \(GL(X)\). Hence \(K=\overline{\pi(G)}\) is compact. Choose
any positive definite Euclidean inner product in the real case, or Hermitian
inner product in the complex case, and average it over \(K\). This gives an
auxiliary inner product \(h\) for which every \(\pi_g\) is unitary.

Define the Riesz identification \(J:X\to X^*\) by
\[
   (Jx)(y)=h(y,x).
\]
In the real case \(J\) is linear; in the complex case it is conjugate-linear.
In both cases it is a bijection and satisfies
\[
   J\pi_g=\bar\pi_gJ,\qquad g\in G.
\]
Indeed, by \(K\)-invariance of \(h\),
\[
   (J\pi_gx)(y)=h(y,\pi_gx)=h(\pi_{g^{-1}}y,x)
      =(\pi_{g^{-1}}^*Jx)(y)=(\bar\pi_gJx)(y).
\]

Apply \(J\) coordinatewise to functions on \(S\), writing
\[
   (\mathcal J f)(s)=J(f(s)).
\]
The intertwining relation shows that \(\mathcal J\) carries
\(C_-^{(1,p)}(G,\pi)\) onto \(C_-^{(1,p^*)}(G,\bar\pi)\). It also commutes
with the \(d\)-operators in the sense that
\[
   d_{\bar\pi}(\mathcal J f)=\mathcal J(d_\pi f),
\]
because each term in \(d_\pi f(s,t)=\pi_s f(s^{-1}t)-f(t)+f(s)\) is sent to
the corresponding term in
\[
   d_{\bar\pi}(\mathcal J f)(s,t)
   =\bar\pi_sJf(s^{-1}t)-Jf(t)+Jf(s).
\]
Therefore \(\mathcal J\) is a bijection from \(\ker d_\pi\) onto
\(\ker d_{\bar\pi}\). In particular, the two kernels have the same finite
dimension.

Now let \(0\ne f\in\ker d_\pi\). Then \(\mathcal J f\in\ker d_{\bar\pi}\), and
the restriction pairing gives
\[
   R_\pi(\mathcal J f)(f)
   =\sum_{s\in S}(Jf(s))(f(s))\,\deg_\omega(s)
   =\sum_{s\in S}h(f(s),f(s))\,\deg_\omega(s)>0.
\]
Thus the range of \(R_\pi\) separates points of \(\ker d_\pi\). Since
\(\dim \ker d_{\bar\pi}=\dim\ker d_\pi=\dim(\ker d_\pi)^*\), this separation
implies that \(R_\pi\) is surjective, hence an isomorphism.
\end{proof}

\section*{Verification notes}

The proof is purely finite-dimensional and does not use either Poincare
constant inequality from Nowak's main theorem. The compactness input is only the
standard fact that the linear isometry group of a finite-dimensional normed
space is compact. The complex case uses the conjugate-linear Riesz map only to
identify the two kernels and separate points; this is enough for complex-linear
surjectivity of \(R_\pi\).

The script
\texttt{code/finite\_dimensional\_kernel\_pairing\_check.py} performs rank
sanity checks for cyclic and free-group-style generating sets with non-orthogonal
conjugates of rotations. It checks that the restriction pairing has full rank in
those examples. These computations are not used in the proof.

\section*{Novelty and limitations}

Cheap run-index searches for arXiv:1107.1896 and phrases around the
finite-dimensional kernel-pairing statement found no duplicate packet. Web
searches on 2026-07-03 for exact phrases involving \(i^*\circ\bar i\),
finite-dimensional isometric representations, and the paper title found no
prior answer beyond the source paper. The result remains a partial answer:
Remark 4.3 is still open here for infinite-dimensional isometric
representations.

\begin{thebibliography}{9}
\bibitem{Nowak2011}
P. W. Nowak,
\emph{Poincare inequalities and rigidity for actions on Banach spaces},
arXiv:1107.1896, 2011.
\end{thebibliography}

\end{document}
